% !PRINTED 687
\weberpage{687}
Nach (9) verschwinden $\alpha_1$, $\beta_1$ in denselben Punkten und mit denselben Ordnungszahlen wie $\alpha$, $\beta$, und daraus ergeben sich, den Darstellungen (2) und (4) entsprechend, die Polygondarstellungen:
\weq{11}{\alpha_1 = \dfrac{\mathfrak{M}}{\mathfrak{A}_1}, \quad \beta_1 = \dfrac{\mathfrak{N}}{\mathfrak{B}_1},}
und aus (5), (7) und (10):
\weq{12}{F_1'(\alpha_1) = \dfrac{\mathfrak{L}}{\mathfrak{A}_1^{n-2}\mathfrak{B}_1^m}.}

Damit ist die Formel (8) \gesperrt{allgemein} bewiesen.

Setzt man nach \S\ 194, (9):
\[
F'(\alpha)\,d\alpha + F'(\beta)\,d\beta = 0,
\]
so folgt aus (8) und \S\ 195, (12):
\weq{13}{F'(\beta) = \dfrac{\mathfrak{R}\mathfrak{Z}_\alpha}{\mathfrak{A}^n\mathfrak{B}^{m-2}}.}

Beide Ableitungen $F'(\alpha)$, $F'(\beta)$ verschwinden also in den Punkten von $\mathfrak{R}$, und $\mathfrak{R}$ hei\ss{}t das Polygon der Doppelpunkte in $(\alpha,\beta)$.

Die Ordnung von $\mathfrak{A}^n\mathfrak{B}^{m-2}$ ist $2n(m-1)$ und daraus ergibt sich f\"ur die Ordnung $2r$ von $\mathfrak{R}$ eine gerade Zahl, n\"amlich:
\[
2r = 2n(m-1) - w_\alpha = 2m(n-1) - w_\beta,
\]
und f\"ur das Geschlecht $p$ ergibt sich nach (2):
\weq{14}{p = (n-1)(m-1) - r.}

\begin{center}
\rule{4cm}{0.4pt}
\end{center}
